\section{Limits at infinity}

Some limits describe behavior near a finite point. Other limits describe behavior far away. When we write
\[
\lim_{x\to\infty}f(x),
\]
we ask what happens to \(f(x)\) as \(x\) becomes larger and larger. When we write
\[
\lim_{x\to-\infty}f(x),
\]
we ask what happens as \(x\) becomes very large in the negative direction.

This is closely related to sequence limits, where we ask what happens as \(n\to\infty\). The difference is that \(n\) runs through natural numbers, while \(x\) runs through real values.

\subsection*{A first example}

Consider
\[
f(x)=\frac1x.
\]
As \(x\) becomes very large and positive, \(1/x\) becomes close to \(0\). Therefore
\[
\lim_{x\to\infty}\frac1x=0.
\]
As \(x\) becomes very large and negative, \(1/x\) also becomes close to \(0\). Therefore
\[
\lim_{x\to-\infty}\frac1x=0.
\]

The sign is different on the two sides, but the values still approach the same number \(0\).

\subsection*{Rational functions}

For rational functions, limits at infinity are often found by comparing the highest powers of \(x\).

\begin{examplebox}
Compute
\[
\lim_{x\to\infty}\frac{3x^2-5x+1}{2x^2+7}.
\]
Divide numerator and denominator by \(x^2\):
\[
\frac{3x^2-5x+1}{2x^2+7}
=
\frac{3-\frac5x+\frac1{x^2}}{2+\frac7{x^2}}.
\]
As \(x\to\infty\), the terms \(1/x\) and \(1/x^2\) tend to \(0\). Hence
\[
\lim_{x\to\infty}\frac{3x^2-5x+1}{2x^2+7}=\frac32.
\]
\end{examplebox}

This is the same dominant-term idea used for sequence limits. The largest powers determine the long-term behavior.

\begin{examplebox}
Compute
\[
\lim_{x\to\infty}\frac{2x+1}{x^2+3}.
\]
The denominator grows faster than the numerator. Divide by \(x^2\):
\[
\frac{2x+1}{x^2+3}=\frac{\frac2x+\frac1{x^2}}{1+\frac3{x^2}}.
\]
The numerator tends to \(0\), and the denominator tends to \(1\). Therefore
\[
\lim_{x\to\infty}\frac{2x+1}{x^2+3}=0.
\]
\end{examplebox}

\subsection*{Horizontal asymptotes}

If
\[
\lim_{x\to\infty}f(x)=L
\]
or
\[
\lim_{x\to-\infty}f(x)=L,
\]
then the graph approaches the horizontal line
\[
y=L
\]
in that direction. This line is called a horizontal asymptote.

\begin{examplebox}
For
\[
f(x)=\frac{3x^2-5x+1}{2x^2+7},
\]
we found
\[
\lim_{x\to\infty}f(x)=\frac32.
\]
Also,
\[
\lim_{x\to-\infty}f(x)=\frac32.
\]
Thus the graph has horizontal asymptote
\[
y=\frac32.
\]
\end{examplebox}

The asymptote is not necessarily a line the graph never touches. It is a line the graph approaches far away.

\subsection*{Infinite limits}

Sometimes a function grows without bound as \(x\to\infty\). For example,
\[
\lim_{x\to\infty}x^2=\infty.
\]
This means the values become larger than any fixed number. It is not convergence to a real number.

Similarly,
\[
\lim_{x\to-\infty}x^3=-\infty.
\]
The values become negative with arbitrarily large magnitude.

\begin{summarybox}
When studying limits at infinity, ask:
\begin{itemize}
\item Is the input going to \(\infty\) or \(-\infty\)?
\item Which terms dominate for large \(|x|\)?
\item Does the function approach a finite number?
\item Does it grow without bound?
\item Is there a horizontal asymptote?
\item Is the behavior different for \(x\to\infty\) and \(x\to-\infty\)?
\end{itemize}
\end{summarybox}

Limits at infinity describe long-term behavior. They tell us what the function does far from the origin.